\chapter*{Exercices de conversion}


\section*{Rappel : Les unités de mesure du débit}
En réseau, on utilise le système décimal pour les débits. Le facteur de conversion est de $1 000$.
\begin{itemize}
    \item 1 kbits/s (kilobit par seconde) = $10^3$ bits/s
    \item 1 Mbits/s (mégabit par seconde) = $10^6$ bits/s ($1\ 000$ kbits/s)
    \item 1 Gbits/s (gigabit par seconde) = $10^9$ bits/s ($1\ 000$ Mbits/s)
\end{itemize}

\subsection*{Le calcul du coût OSPF}
Contrairement au RIP qui compte simplement le nombre de sauts, l'OSPF calcule un coût basé sur la bande passante. La formule est la suivante :

$$Coût = \frac{Debit\ de\ reference}{Debit\ reel}$$

\begin{Comprendre}
    Règle d'or : Pour que la division soit valide, le \textbf{débit de référence} et le \textbf{débit réel} doivent être exprimés dans la \textbf{même unité}.
\end{Comprendre}


\section*{Exercice 1 : conversion}
Convertissez les valeurs suivantes pour préparer vos calculs de coût :
\begin{enumerate}
    \item Liaison Ethernet standard : $10\ 000\ 000$ bits/s = ............ Mbits/s
    \item Liaison Fast Ethernet : $100$ Mbits/s = ............ Gbits/s
    \item Liaison Fibre : 10 Gbits/s = ............ Mbits/s
    \item Liaison ADSL : $20$ Mbits/s = ............ bits/s
\end{enumerate}



\section*{Exercice 2 : La pyramide des débits}
Complétez les égalités suivantes en utilisant les puissances de 10 (sous la forme $10^n$).
\begin{itemize}
    \item 1 kbits/s = $10^3$ bits/s
    \item  1 Mbits/s = $10^{...}$ bits/s
    \item 1 Gbits/s = $10^{...}$ bits/s
    \item 1 Gbits/s = $10^{...}$ Mbits/s
    \item 100 Mbits/s = $10^{...}$ bits/s
    \item 2 Gbits/s = $2 \times 10^{...}$ bits/s
\end{itemize}
 En résumé, calculez le coût pour les interfaces suivantes :

\begin{table}[!ht]
    \centering
    \begin{tabular}{|l|l|l|l|}
    \hline
        Type de liaison & Débit réel & Calcul (avec unités converties) & Coût (arrondi à l'entier sup) \\ \hline
        \textbf{Ethernet} & 10 Mbits/s & $\frac{100\ Mbits/s}{10\ Mbits/s}$ & $10$\\ \hline
        \textbf{Fast Ethernet} & 100 Mbits/s & ............................................................ & ............ \\ \hline
        \textbf{Giga Ethernet} & 1 Gbits/s & ............................................................ & ............ \\ \hline
        \textbf{Fibre} & 10 Gbits/s & ............................................................ & ............ \\ \hline
    \end{tabular}
\end{table}


\section*{Exercice 3 : Simplification de fractions (Le réflexe OSPF)}
Pour calculer le coût OSPF, on divise souvent deux puissances de 10. Rappel mathématique : $\frac{10^a}{10^b} = 10^{a-b}$.

Calculez le résultat des divisions suivantes (sans calculatrice) :
\begin{enumerate}
    \item $\frac{10^8}{10^7} = 10^{8-7} = 10^1 =10$
    \item  $\frac{10^8}{10^8} = ............ =............$
    \item  $\frac{10^8}{10^6} = ............ = ............$
    \item  $\frac{10^8}{10^9} = ............ = ............$ 
\end{enumerate}
 



\section*{Exercice 4 : Mise en situation "Unités dépareillées"}
Pour chaque situation, convertissez le débit du lien dans la même unité que le débit de référence ($10^8$ bits/s), puis calculez le coût.
\subsection*{Liaison Satellite :}
\begin{itemize}
    \item Débit de référence : $10^8$ bits/s
    \item Débit réel : $50$ Mbits/s
    \item Conversion du réel en bits/s :..............................
    \item Calcul du coût :..................................................
\end{itemize}


\subsection*{Liaison Backbone :}
\begin{itemize}
    \item Débit de référence : $10^8$ bits/s
    \item Débit réel : $10$ Gbits/s
    \item Conversion du réel en bits/s : ..............................
    \item Calcul du coût : ..................................................

\end{itemize}



\section*{Exercice 5 : Défi inverse}
Le coût OSPF d'une liaison est de $100$. Le débit de référence est configuré à $10^8$ bits/s.
\begin{enumerate}
    \item Trouvez le débit réel de la liaison en $bits/s$.
    \item Convertissez ce résultat en $Mbits/s$.
\end{enumerate}
 
\section*{Exercice 6 : Analyse et réflexion}
\begin{enumerate}
\item \textit{Problème de cohérence :} Si un élève divise un débit de référence de $100$ Mbits/s par un débit réel de $1$ Gbits/s sans convertir, il obtient un coût de $100$. Est-ce logique ? Pourquoi ?
\item \textit{Comparaison RIP :} Si un paquet doit traverser 3 routeurs reliés en Giga Ethernet (1 Gbits/s) ou 2 routeurs reliés en Ethernet (10 Mbits/s) :
    \begin{enumerate}
        \item Quel chemin choisira RIP ?
        \item Quel chemin choisira OSPF ?
    \end{enumerate}
    
\end{enumerate}



\begin{Comprendre}
\textbf{Pourquoi insister sur les puissances de 10 ?}\\
Dans un réseau, le coût est une valeur \textit{sans unité}. 
Si vous divisez $100$ Mbits/s par $1$ \textit{Gbits/s}, vous obtenez $100$. C'est faux car vous avez divisé des "pommes" par des "poires". \\
En utilisant les puissances de 10, on s'assure de l'homogénéité : 
$$\frac{10^8 \text{ bits/s}}{10^9 \text{ bits/s}} = 10^{-1} = 0,1$$
\end{Comprendre}


\section*{Exercice 7 : Retrouver le débit à partir du coût}

Dans cet exercice, on considère que le \textit{débit de référence} du réseau est configuré à $10^8$ bits/s (soit 100 Mbits/s).

La formule du coût est : $Coût = \frac{Débit\ de\ référence}{Débit\ réel}$\\

Pour chaque cas ci-dessous, retrouvez le débit réel de la liaison. Donnez le résultat en \textit{puissances de 10}, puis convertissez-le dans l'unité la plus adaptée (Mbits/s ou Gbits/s).

\subsection*{Cas n°1 : La liaison lente}
Le routeur affiche un coût de $100$ pour une interface.
\begin{enumerate}
    \item Écrivez l'équation : $100 = \frac{10^8}{Débit\ réel}$
    \item Isolez le débit réel : $Débit\ réel = \frac{10^8}{10^2} = ............$ bits/s
    \item Convertissez en Mbits/s : ...........................
\end{enumerate}


\subsection*{Cas n°2 : La liaison standard}
Le coût calculé par OSPF est de $10$.
\begin{enumerate}
    \item Trouvez le débit réel en bits/s en utilisant les puissances de 10 : ............................
    \item Convertissez le résultat en Mbits/s : ............................
\end{enumerate}

\subsection*{Cas n°3 : La liaison "référence"}
Le coût affiché est de $1$.
\begin{enumerate}
    \item Quelle est la valeur du débit réel en bits/s ? ............................
    \item À quelle unité de base cela correspond-il (en Mbits/s) ? ............................
\end{enumerate}

\subsection*{Cas n°4 : Le piège du débit de référence (Défi)}
L'administrateur change le débit de référence pour l'ajuster aux nouvelles fibres optiques. Le nouveau débit de référence est $10^9$ bits/s (1 Gbits/s).
Une interface affiche un coût de $50$.
\begin{enumerate}
    \item  Calculez le débit réel en bits/s : ............................
    \item Convertissez ce débit en Mbits/s : ............................
\end{enumerate}

\begin{Comprendre}
    Pour trouver le débit, on utilise la transformation de la formule :
$$Debit\ réel = \frac{Debit\ de\ reference}{Cout}$$
\end{Comprendre}
